<!-- CS 452


Lecture 4:

"Fake a stack"

- Register File
   - SP and PC are not part of register file (X0, ..., x30, x31
   - SP for each EL (exception level): 
      - SP_EL0 (USER),
      - SP_EL1 (KERNEL)
      - SP_EL2 (HYPERVISOR) -> boot.s
      - SP_EL3 (SECURE MONITOR)
   - System registers 
      - MRS (read) or MSR (write)
      - DO CONSTANTLY SET SPSel to 0 (should be set to 1)
      - ELR_EL1 (Exception Link Register) -- where to return after exception
      - ESR_EL1 (Exception Syndrome Register) -- tells why the exception occured
      - SPSR_EL1 -- saved processor state register

      - Exception Vector
         - VBAR_EL1 (Vector Base Address Register) -- where to find the exception handler 
         - Table 4 in exception document 
            - Want lowerEL 64
               - Syncronous Exception (used for SVC)
               - IRQ (Interrupt Request)
               - FIQ (Fast Interrupt Request)
               - Error 

               - each section has 128 bytes ~ 32 instructions
            - Recommended to set up all 4 with dummy routines

      

- System Call
   - SVC N: "branch to handler" (handler is essentially a subroutine)
   - Handler:
      - save arithmetic register 
      - (Save ELR/SPSR)
      - perform service
      - restore registers
      - (restore ELR/SPSR)
      - eret (exception return) -- return from exception (use this only as stepping stone first)
   - Execution mode increments by 1 (EL0 -> EL1)
   - Sets ___EL1 registers
   - N gets put in lower 16 bits of ESR_EL1


- Processor State (pstate -- psr)
   - execution state
   - execution level
   - interrupt flags
   - condition codes:
      - N (negative)
      - Z (zero)
      - C (carry)
      - V (overflow)
-->